% Part: proof-theory
% Chapter: normalization
% Section: introduction

\documentclass[../../../include/open-logic-section]{subfiles}

\begin{document}

\olfileid{pt}{nor}{seg}

\olsection{ખંડો અને કટ}

!!^a{introduction} નિયમ
પછી !!a{elimination}
નિયમ આવે તે
!!a{proof}માં ચકરાવો
છે; !!a{proof}નું
સામાન્યીકરણ એટલે
એવા ચકરાવા
દૂર કરવા.
જોકે~\Elim{\lor}
અને~\Elim{\lexists}
હોવાને કારણે
``દૂર કરવાનો
ચકરાવો'' શું છે
તેની વ્યાખ્યા
પણ થોડી જટિલ
બને છે. સમસ્યા
એ છે કે
આ નિયમોને લીધે
ચકરાવામાં
!!{elimination}
નિયમ !!{introduction}
નિયમની \emph{તરત}
પછી આવે જ
તે જરૂરી નથી.
વચ્ચે~\Elim{\lor}
કે~\Elim{\lexists}
આવી શકે છે,
જેમ કે:
\[
\AxiomC{}
\DeduceC{$\lexists[x][!C(x)]$}
\AxiomC{$\Discharge{!A}{x}$}
\DeduceC{$!B$}
\DischargeRule{\Intro{\lif}}{x}
\UnaryInfC{$!A \lif !B$}
\DischargeRule{\Elim{\lexists}}{y}
\BinaryInfC{$!A \lif !B$}
\AxiomC{}
\DeduceC{$!A$}
\RightLabel{\Elim{\lif}}
\BinaryInfC{$!B$}
\DisplayProof
\]
ખરેખર, ગમે તેટલા
\Elim{\lor} અથવા
\Elim{\lexists}
નિયમો~\Intro{\lif}
અને~\Elim{\lif}
વચ્ચે આવી શકે
છે, જેમ કે:
\[
\AxiomC{}
\DeduceC{$\lexists[x][!C(x)]$}
\AxiomC{}
\DeduceC{$!D \lor !E$}
\AxiomC{$\Discharge{!A}{x}$}
\DeduceC{$!B$}
\DischargeRule{\Intro{\lif}}{x}
\UnaryInfC{$!A \lif !B$}
\AxiomC{$\Discharge{!A}{x}$}
\DeduceC{$!B$}
\DischargeRule{\Intro{\lif}}{x}
\UnaryInfC{$!A \lif !B$}
\RightLabel{\Elim{\lor}}
\TrinaryInfC{$!A \lif !B$}
\DischargeRule{\Elim{\lexists}}{y}
\BinaryInfC{$!A \lif !B$}
\AxiomC{}
\DeduceC{$!A$}
\RightLabel{\Elim{\lif}}
\BinaryInfC{$!B$}
\DisplayProof
\]
આ દાખલો બતાવે
છે કે એક જ
\Elim{\lif} અનેક
ચકરાવામાં સામેલ
હોઈ શકે છે,
કારણ કે તેના
પહેલાં એકથી વધુ
\Intro{\lif} આવે
છે. આ બધી
શક્યતાઓ સમાવવા
\Log{N1i}-સાબિતીમાં
!!{formula} આવૃત્તિઓના
અમુક પ્રકારના
ક્રમ, જેને
\emph{ખંડો} કહીએ,
તે વડે ચકરાવા
વ્યાખ્યાયિત કરીએ.
આપણા દાખલામાં
તે~$!A \lif !B$ની
આવૃત્તિઓનો એવો
ક્રમ છે જેમાં
$\Elim{\lor}$ કે
$\Elim{\lexists}$ના
ગૌણ આધારવિધાન
પછી એ નિયમનું
નિષ્કર્ષ આવે
છે. અહીં આવા
બે ખંડ છે,
અને દરેકમાં~$!A \lif !B$ની
ત્રણ આવૃત્તિઓ
છે. હવે ઔપચારિક
વ્યાખ્યા આપીએ.

\begin{defn}
\Log{N1i}ની !!{proof}~$\delta$માં
એક જ !!{formula}~$!A$ની
આવૃત્તિઓનો ક્રમ
$!A_1$, \dots, $!A_n$
એ જ~$\delta$માં
નીચેની શરતો પૂરી
કરે તો તે
\emph{ખંડ} છે:
\begin{enumerate}
  \item $!A_1$ એ
  \Elim{\lor} કે
  \Elim{\lexists}નું
  નિષ્કર્ષ ન હોય;
  \item $!A_n$ એ
  \Elim{\lor} કે
  \Elim{\lexists}નું
  ગૌણ આધારવિધાન
  ન હોય;
  \item $i = 1$,
  \dots,~$n-1$
  માટે~$!A_i$ એ
  \Elim{\lor}
  અથવા~\Elim{\lexists}નું
  ગૌણ આધારવિધાન
  હોય અને
  $!A_{i+1}$
  તે જ અનુમાનનું
  નિષ્કર્ષ હોય.
\end{enumerate}
\sourcecorrection{OLMD-259}{મૂળની ત્રીજી શરતમાં પહેલાની આવૃત્તિ $!A_i$ પછીની આવૃત્તિ $A_{i+1}$ લખાઈ છે; બંને એક જ $!A$ સૂત્રની આવૃત્તિઓ હોવાથી બીજામાં પણ $!$ જરૂરી છે.}
તે ખંડની
\emph{લંબાઈ}~$n$ છે.
\end{defn}

દરેક ખંડ દૂર
કરવા યોગ્ય
ચકરાવો નથી.
જે એવા છે
તેમને \emph{કટ}
કહીએ છીએ.

\begin{defn}
\emph{કટ} એવો
ખંડ છે જેમાં
$!A_n$ એ
!!a{elimination}
અનુમાનનું મુખ્ય
આધારવિધાન હોય
અને કાં તો
$n=1$ હોય તથા
$!A_1 = !A_n$
એ !!a{introduction}
અનુમાનનું અથવા
\FalseInt{}નું
નિષ્કર્ષ હોય,
અથવા~$n>1$
હોય.
\end{defn}

નોંધો કે કટ-ખંડની
શરૂઆત !!a{introduction}
નિયમના નિષ્કર્ષથી
થવી જરૂરી નથી.
તેનો અંત
!!a{elimination}
નિયમના મુખ્ય
આધારવિધાનમાં
થવો પૂરતો છે.
જોકે લંબાઈ~$1$નો
ખંડ કટ ગણાય
તે માટે તે
!!{introduction}
નિયમનું અથવા
\FalseInt{}નું
નિષ્કર્ષ હોવું
જોઈએ.
\sourcecorrection{OLMD-260}{મૂળનો સારાંશ લંબાઈ એકના કિસ્સામાં ફક્ત પરિચય-નિયમ કહે છે, જ્યારે તરત પહેલાંની કટની વ્યાખ્યામાં \FalseInt{} પણ સ્પષ્ટ સામેલ છે.}

\begin{defn}
કટની \emph{કક્ષા}~$\depth{!A}$
છે. !!a{proof}~$\delta$ની
\emph{કટ-કક્ષા}~$\cutrank{\delta}$
એ~$\delta$માંના
કટોની મહત્તમ
કક્ષા છે.
કટની કક્ષા
$\cutrank{\delta}$
હોય, એટલે
મહત્તમ કક્ષા
હોય, તો તે
$\delta$માં
\emph{મહત્તમ}
કટ છે.
$\delta$ની
\emph{કટ-લંબાઈ}
એ તેના બધા
મહત્તમ કટોની
લંબાઈનો સરવાળો
છે.
\end{defn}

\end{document}
